% ICLR paper fragment: descriptive analysis of the synthetic observations.
% Required packages in the parent document: graphicx, booktabs, amsmath.
% The figure directory may be redefined before \input{...} when the paper is
% compiled outside the repository root.
\providecommand{\syntheticdatafigdir}{figures/synthetic_appendix}

\subsection{Synthetic-data diagnostics: repeatability, cross-view agreement,
and latent-profile fidelity}
\label{sec:synthetic_data_diagnostics}

The synthetic collection is designed to separate a transcript-specific signal
from observation effects that are reproducible within a dataset.  All
conditions share the same programmed kinetic profile $K_t$, but differ through
ten sequence-dependent observation rules: four $5'$ dinucleotide biases, four
$3'$ dinucleotide biases, and AU-rich and GC-rich sequence-context biases.  For
each rule, we sample two count replicas at nominal depths of $0.25$, $2$, and
$20$ reads per codon.  The analysis below covers 19,283 aligned transcripts and
all modeled sense codons; the appended terminal position is excluded.  We do
not include the separately sampled unbiased observation condition.  In
particular, $K_t$ denotes the programmed kinetic target and not that sampled
condition.  It is also not an exact noiseless expected count profile, because
the simulator contains traffic and count-sampling effects in addition to the
programmed dwell-time profile.  We therefore use the matching mean-one TASEP
occupancy $q_t^{(r)}$ as the primary pre-bias target for sampled replica $r$,
and $\bar q_t=(q_t^{(1)}+q_t^{(2)})/2$ for the arithmetic replica consensus.
Agreement with $K_t$ is retained separately as a stricter measure of
programmed-kinetic fidelity.  We write $Y^{(r)}_{dt}$ for the sampled
position-wise profile of replica $r$ from dataset condition $d$ and transcript
$t$.

We organize the diagnostic analysis around four expectations.  These are
descriptive expectations, rather than preregistered statistical hypotheses.
\textbf{H1 (sampling precision):} within-dataset replica agreement should
increase with read depth.  \textbf{H2 (shared signal):} different bias
conditions should retain positive agreement because they share $K_t$; however,
agreement need not approach one because their observation effects differ.
Two forces act on the cross-dataset correlation as depth increases: reduced
sampling noise reveals the common signal, whereas improved precision also
reveals systematic bias more clearly.  Its net direction is therefore an
empirical question.  \textbf{H3 (replica averaging):} the arithmetic mean of
two replicas should agree more closely with $\bar q_t$ than either replica
agrees with its matching $q_t^{(r)}$, with the largest gain at low depth.
\textbf{H4 (repeatability is not
fidelity):} a systematic bias may be highly reproducible without being close
to the pre-bias occupancy.  Consequently, replica agreement is a measure of technical
repeatability, not a sufficient measure of kinetic fidelity.

\begin{table*}[t]
    \centering
    \footnotesize
    \caption{Headline synthetic-data diagnostics, reported as mean $\pm$ the
    half-width of a two-sided 95\% confidence interval.  Each input value is a
    transcript-level PCC median.  Within-dataset and TASEP-target summaries use a
    Student-$t$ interval across the ten programmed bias conditions.
    Cross-dataset estimates average the 45 pair medians, with uncertainty from
    a leave-one-dataset-out jackknife over the ten datasets so that shared
    members do not make the pairs appear independent.  ``Single replica''
    first averages the two replica-specific medians within each bias.  These
    intervals describe variation across the selected synthetic conditions,
    not population-level biological uncertainty.}
    \label{tab:synthetic_data_diagnostics}
    \setlength{\tabcolsep}{4pt}
    \begin{tabular}{@{}lcccc@{}}
        \toprule
        \shortstack{Reads/\\codon}
        & \shortstack{Replica\\agreement}
        & \shortstack{Cross-dataset\\agreement}
        & \shortstack{Single replica\\vs.\ matching $q_t^{(r)}$}
        & \shortstack{Replica mean\\vs.\ $\bar q_t$} \\
        \midrule
        $0.25$ & $0.186 \pm 0.042$ & $0.190 \pm 0.039$ & $0.257 \pm 0.005$ & $0.322 \pm 0.011$ \\
        $2$    & $0.538 \pm 0.067$ & $0.266 \pm 0.064$ & $0.450 \pm 0.035$ & $0.499 \pm 0.052$ \\
        $20$   & $0.731 \pm 0.040$ & $0.324 \pm 0.098$ & $0.537 \pm 0.070$ & $0.559 \pm 0.082$ \\
        \bottomrule
    \end{tabular}
\end{table*}

\paragraph{Read depth controls technical repeatability.}
Figure~\ref{fig:synthetic_replica_agreement} compares the full transcript-level
distribution of
$\operatorname{PCC}(Y^{(1)}_{dt},Y^{(2)}_{dt})$ for every bias condition.  The
median across the ten condition medians rises from $0.170$ at $0.25$ reads per
codon, to $0.534$ at $2$ reads per codon, and to $0.735$ at $20$ reads per
codon.  The increase occurs for every bias rule, supporting H1.  The remaining
between-condition spread is substantial even at the highest depth
($0.628$--$0.791$).  The AU-rich and GC-rich conditions additionally show broad
transcript distributions because the corresponding rules are sparse and
sequence dependent.  For example, the AU-rich rule selects no codon in
$47.6\%$ of transcripts and has a median selected-codon fraction of only
$0.22\%$.  Thus the box width and location jointly reflect count precision,
bias prevalence, and the magnitude of the resulting positional peaks.

\begin{figure*}[t]
    \centering
    \includegraphics[width=\textwidth]{\syntheticdatafigdir/replica_agreement_boxplots_by_dataset_and_depth.pdf}
    \caption{\textbf{Within-dataset replica agreement across read depths.}
    For each programmed bias, adjacent boxes show the transcript-level
    distribution of $\operatorname{PCC}(Y^{(1)}_{dt},Y^{(2)}_{dt})$ at the
    three nominal read depths.  Boxes show the interquartile range and median;
    whiskers show the 5th and 95th percentiles.  Each box contains 19,283
    transcripts.  The monotonic depth effect measures increasing technical
    repeatability and does not by itself imply recovery of $K_t$.}
    \label{fig:synthetic_replica_agreement}
\end{figure*}

\paragraph{Different observation rules share signal but are not
interchangeable.}
For every transcript, we form the arithmetic two-replica consensus in each
dataset and correlate it with the consensus from every other dataset.  Each
cell in Figure~\ref{fig:synthetic_cross_dataset_agreement} is the median of
19,283 transcript-level PCCs; it is not a pooled codon-level correlation.  The
median across the 45 off-diagonal pair medians increases from $0.185$ to
$0.253$ and $0.301$ with depth.  This direction suggests that, in these data,
removing count noise reveals the shared component more strongly than it reveals
between-rule differences.  Nevertheless, high-depth correlations remain
heterogeneous ($0.164$--$0.586$), so the datasets cannot be treated as repeated
measurements of one common observation profile.  At all depths, the largest
median is obtained for the $3'$-AA/AU-rich pair and the smallest for the
$3'$-CC/$3'$-UU pair.  These identities are useful descriptions of this
realization, but should not be elevated to general properties of those
sequence motifs.

\begin{figure*}[t]
    \centering
    \includegraphics[width=\textwidth]{\syntheticdatafigdir/cross_dataset_consensus_correlation_by_depth.pdf}
    \caption{\textbf{Cross-dataset agreement of the arithmetic replica
    consensuses.}  Each off-diagonal entry is the median, over transcripts, of
    the PCC between two bias-specific consensus profiles.  The diagonal is one
    by definition, and all depths use the same color scale.  Agreement rises
    with read depth but remains heterogeneous and far from unity, revealing
    persistent view-specific structure.}
    \label{fig:synthetic_cross_dataset_agreement}
\end{figure*}

The cross-dataset matrices require an additional construction caveat.  The
bias conditions reuse the same two TASEP occupancy trajectories, so they are
paired synthetic views rather than independently regenerated experiments.
Positions unaffected by either rule can retain exactly equal
sampled counts, particularly at low depth where zeros are common.  For
example, approximately $5\%$ of AU-rich/GC-rich transcript profiles are
exactly identical in each depth-specific realization.  Cross-dataset PCC
therefore quantifies similarity among the provided paired files; it should not
be interpreted as the agreement expected between independent experimental
replicates or independent simulator seeds.

\paragraph{Averaging reduces sampling noise but cannot remove systematic
bias.}
Figure~\ref{fig:synthetic_kinetic_agreement} compares each replica and their
arithmetic mean with the matched pre-bias occupancy targets.  Specifically,
$Y_{dt}^{(r)}$ is compared with $q_t^{(r)}$, while
$(Y_{dt}^{(1)}+Y_{dt}^{(2)})/2$ is compared with $\bar q_t$.  The two individual
replicas have nearly identical distributions within every condition, as
expected from exchangeable sampling.  Before technical bias and count
sampling, the two finite TASEP trajectories have median transcript-level PCC
$0.866$ (interquartile range $0.849$--$0.882$), demonstrating that trajectory
recording itself contributes a smaller but nonzero source of replicate
variation.  Across bias conditions, the median
consensus--$\bar q_t$ PCC is $0.323$, $0.499$, and $0.549$ at the three depths,
compared with $0.256$, $0.452$, and $0.530$ for a typical individually matched
replica.  The corresponding median gains from averaging are $0.067$, $0.046$,
and $0.018$, consistent with H3 and with diminishing returns once sampling
noise is reduced.

\begin{figure*}[t]
    \centering
    \includegraphics[width=\textwidth]{\syntheticdatafigdir/tasep_occupancy_agreement_boxplots_by_dataset_and_depth.pdf}
    \caption{\textbf{Agreement with matched pre-bias TASEP occupancy.}
    Blue and orange show transcript-level PCCs between sampled replica $r$ and
    its own normalized occupancy $q_t^{(r)}$; green compares the arithmetic
    sampled profile with $\bar q_t$.  Valid codon positions are identical in
    all comparisons.  Averaging improves agreement most at low depth, whereas
    persistent differences at high depth expose systematic sequence-dependent
    observation effects.}
    \label{fig:synthetic_kinetic_agreement}
\end{figure*}

For comparison, the corresponding median consensus--$K_t$ PCCs are $0.310$,
$0.478$, and $0.527$.  Their modest deficits relative to the occupancy-based
values quantify the additional separation introduced by traffic and finite
trajectory recording.  Thus $q_t^{(r)}$ is the appropriate proximal target
for auditing count generation, whereas $K_t$ remains the appropriate target
for evaluating recovery of programmed sequence-dependent kinetics.

Crucially, the most repeatable condition is not necessarily the most faithful.
At $20$ reads per codon, the $3'$-CC condition has the highest median replica
agreement ($0.791$) but the lowest median consensus--$\bar q_t$ agreement ($0.442$).
Conversely, the sparse AU-rich rule has lower replica agreement ($0.628$) but
the highest consensus--$\bar q_t$ agreement ($0.788$), because it leaves most
positions and many complete transcripts unaffected.  Across the ten
condition-level medians, the Spearman association between replica agreement
and consensus--$\bar q_t$ agreement is $-0.903$, $-0.976$, and $-0.988$ at increasing
depth.  This negative association is descriptive of the ten deliberately
constructed rules and one simulator/count seed; it is not a population-level
estimate and we do not attach a codon-level significance test to it.  It
nevertheless provides a direct synthetic demonstration of H4: greater
repeatability can indicate a more consistently expressed observation bias.

\paragraph{Implication for the learning problem.}
Taken together, the diagnostics establish the intended stress test.  The
datasets contain a common signal that becomes clearer with depth, yet retain
substantial and reproducible dataset-specific structure.  Replica averaging
attenuates stochastic noise but does not recover the pre-bias occupancy in the
presence of systematic bias.  A model evaluated only on observation fit
or within-dataset repeatability can therefore favor a biased solution.  This
motivates learning a shared profile from multiple heterogeneous views while
representing their observation effects separately.  These diagnostics validate
the structure of the synthetic challenge; they do not, by themselves,
establish identifiability of the learned decomposition or biological realism
of the programmed bias rules.

% Generated numerical sources:
% analyses/artifacts/synthetic/individual_dataset/replica_agreement_summary.csv
% analyses/artifacts/synthetic/individual_dataset/cross_dataset_correlation_summary.csv
% analyses/artifacts/synthetic/individual_dataset/kinetic_target_agreement_summary.csv
% analyses/artifacts/synthetic/individual_dataset/bias_prevalence_summary.csv
% analyses/artifacts/synthetic/individual_dataset/reproducibility_fidelity_association.csv
% analyses/artifacts/synthetic/individual_dataset/headline_mean_ci_summary.csv
% analyses/artifacts/synthetic/individual_dataset/tasep_occupancy_agreement_summary.csv
% analyses/artifacts/synthetic/individual_dataset/headline_tasep_mean_ci_summary.csv
% analyses/artifacts/synthetic/individual_dataset/tasep_occupancy_replica_agreement_summary.csv
